\chapter{Evaluatie en persoonlijke reflectie}
\label{chap:reflectie}

\section{Evaluatie van het stageplan}
Het oorspronkelijke plan legde terecht de nadruk op dataverzameling, lokale opslag, visualisatie en voorspelling. De precieze technische invulling veranderde echter sterk. SQLite werd vervangen door transparante JSONL-, JSON- en CSV-bestanden. Replay werd uit de eindapp verwijderd en vervangen door een testgerichte simulator. Een complex AI/ML-traject maakte plaats voor een uitlegbare sectorvoorspelling. Tegelijk groeide de race-engineeringfunctionaliteit met meerdere sessiemodi, pitstopplanning, providerondersteuning, meerdere dashboards, grafieken en automatische releases.

Die verschuiving was verantwoord. Binnen zes weken was het belangrijker om de volledige live keten betrouwbaar te maken dan om een complex model bovenop onzekere data te bouwen. De moeilijkste problemen zaten niet in de rekenkundige formules, maar in de betekenis en kwaliteit van timingvelden. Zonder correcte providerinterpretatie zou een geavanceerd model alleen preciezere foutieve resultaten produceren.

\section{Zelfstandig werken}
De technische uitvoering gebeurde grotendeels zelfstandig. Dit vereiste dat problemen systematisch werden verkleind: eerst de ruwe snapshot inspecteren, dan de parseroutput, vervolgens het opgeslagen laprecord en ten slotte de analyse. Door functies uit de renderer te halen konden fouten lokaal met kleine inputs worden gereproduceerd.

Zelfstandigheid betekende ook beslissen wanneer een vraag technisch kon worden opgelost en wanneer domeinfeedback nodig was. De betekenis van een referentietijdregel of geldige pitstop kan niet uit code worden afgeleid. In zulke gevallen was overleg met de opdrachtgever noodzakelijk voordat verdere implementatie zinvol was.

\section{Communicatie en requirementsanalyse}
Het frequente overleg met de klant was een centraal deel van de stage. Schermafbeeldingen en concrete scenario's bleken effectiever dan abstracte technische uitleg. Een vraag als ``waar staan we na een pitstop van 75 seconden?'' maakte onmiddellijk zichtbaar welke data en aannames ontbraken. De implementatie kon vervolgens worden uitgelegd in termen van opeenvolgende verschillen, rondeachterstanden en klasseposities.

Ik leerde requirements niet als vaste lijst te behandelen. De opdrachtgever kon pas na gebruik beoordelen of een veld groot genoeg was, of een kleur voldoende opviel en of een vergelijking tijdens een race werkelijk nuttig was. Tegelijk moest ik bewaken dat elke wijziging een duidelijke betekenis hield en niet alleen visueel plausibel werd.

\section{Software-engineeringvaardigheden}
De stage bracht verschillende onderdelen uit de opleiding samen: modulariteit, datastructuren, foutafhandeling, testontwerp, user-interfaceontwerp en distributie. Vooral separation of concerns werd praktisch belangrijk. Pure shared modules bleken niet alleen beter testbaar, maar maakten ook meerdere dashboards en grafiekvensters mogelijk zonder logica te kopiëren.

Het project maakte ook het verschil duidelijk tussen syntax en semantiek. Een parser kan een tekst perfect lezen en toch een fout resultaat produceren wanneer \code{5L} als vijf seconden wordt behandeld. Betrouwbare software vereist daarom domeinconstraints en tests met betekenisvolle voorbeelden.

\section{Projectmatig werken}
De zes weken dwongen tot prioritering. Nieuwe ideeën, zoals weersafhankelijke strategie, bandenkeuze en een uitgebreid FCY-model, waren relevant maar konden niet allemaal volledig worden uitgewerkt. De gekozen werkwijze was telkens eerst een minimale end-to-endfunctie bouwen, die met realistische data testen en pas daarna verfijnen.

Versiebeheer en releases vormden een nuttige afsluiting van iteraties. Featurebranches werden via pull requests naar main gebracht. Tags genereerden reproduceerbare platformbuilds. Automatische tests verminderden het risico dat een nieuwe pitstop- of parserwijziging oudere functionaliteit brak.

\section{Wat beter kon}
Een formele provider-specificatie aan het begin had sommige interpretatiefouten kunnen voorkomen. In de praktijk waren de feeds echter vooral via demo's en screenshots beschikbaar. Ook had een vroeg centraal model voor ``geldige tempo-data'' duplicatie kunnen vermijden; dit inzicht ontstond pas nadat neutralisatieronden op meerdere plaatsen fout werden meegenomen.

De README groeide mee met de ontwikkeling en bevat daardoor op enkele plaatsen historische formuleringen. Voor een volgende fase is het zinvol documentatie per release te controleren tegen de actuele setupflow. Ook digitale signing, notarization en een formeel migratiepad voor opgeslagen data zijn nodig voordat de app breder buiten de huidige testcontext wordt verspreid.

\section{Meerwaarde}
Voor de opdrachtgever levert de stage een werkende basis die ruwe timing vertaalt naar teamgerichte inzichten en verder kan worden uitgebreid. Voor mij lag de grootste meerwaarde in het bouwen van software waarvan een fout tijdens live gebruik directe operationele gevolgen kan hebben. Dat maakte voorzichtigheid, uitlegbaarheid en testen concreter dan in een geïsoleerde academische opdracht.

